\chapter*{Önsöz}
\addcontentsline{toc}{chapter}{Önsöz}
\markboth{Önsöz}{Önsöz}

Mantık, iki bin yılı aşkın bir süre boyunca kâğıt ve kalemle yapılan bir
disiplindi. Aristoteles'in kıyasları, Frege'nin \emph{Begriffsschrift}'i,
Gentzen'in doğal çıkarım sistemi ve Gödel'in eksiklik teoremleri hep bu
gelenekte doğdu. Yirmi birinci yüzyılda ise mantığın kendisi değişmese de onu
\emph{yapma biçimimiz} değişti: Bugün bir ispatı bilgisayara denetletebiliyor,
mantık sistemlerini birer veri yapısı olarak tanımlayıp onlar hakkındaki
teoremleri makinede doğrulayabiliyoruz.

Bu kitap, bu dönüşümü Türkçe okura sıfırdan anlatmak amacıyla yazıldı. Araç
olarak \Lean{} ispat asistanını kullanıyoruz \citep{demoura2021lean4}. \Lean{}
hem bir programlama dili hem de bir ispat denetleyicisidir. İçine yazdığınız
her ispatı, küçük ve güvenilir bir çekirdek adım adım denetler. Bir ispat
\Lean{}'den geçtiyse, o ispatta mantıksal bir boşluk kalmamış demektir.

\section*{Kitabın iki yüzü}

Kitap iki farklı okur kitlesini birlikte gözetiyor:

\begin{description}[font=\baslikfont\bfseries\color{anarenk}]
  \item[Ders kitabı olarak] Kısım~I ve~II, matematik, bilgisayar mühendisliği
  ve felsefe bölümlerinde bir dönemlik ``Mantık ve İspat'' dersine karşılık
  gelir. Önceden programlama bilgisi gerekmez. Her kavram önce kâğıt üzerinde,
  ardından \Lean{}'de, en sonunda da bir görselle anlatılır.
  \item[El kitabı olarak] Kısım~III ve~IV, metamantık (geçerlilik, tamlık,
  kesme eleme) ve tip kuramı ($\lambda$-hesabı, bağımlı tipler, \Lean{}'in
  aksiyomları) konularında araştırmacıya başvuru kaynağı olmayı hedefler.
  Bu bölümlerdeki \emph{Araştırma Notu} kutuları güncel literatüre ve açık
  problemlere yönlendirir. Kısım~V'teki on iki proje ise okuru kendi
  formalizasyon çalışmasına hazırlar.
\end{description}

\section*{İki cilt}

Elinizdeki kitap, iki ciltlik bir dizinin ilk cildidir. \textbf{Cilt~I:
Temeller}, sıfırdan başlayıp önermeler mantığının metateorisine ve
\Lean{}'in kendi tip kuramına kadar ilerler; kendi başına tamamlanmış bir
bütündür. \textbf{Cilt~II: İleri Konular} ise birinci mertebe mantığın
formalizasyonu, kip mantığı, Gödel'in eksiklik teoremleri, küme kuramı,
otomasyon, metaprogramlama ve yapay zekâ destekli teorem ispatını ele
alacaktır. Cilt~I'de ileriye dönük göndermeler ``Cilt~II, Bölüm~$n$''
biçiminde yazılmıştır; Cilt~II'nin planı, Cilt~I'in sonundaki
\emph{Sonsöz}'de yer alır.

\section*{Kitabın omurgası}

Kitabı bir arada tutan tek bir fikir vardır: \textbf{Curry--Howard eşlemesi}.
Bir önerme bir \emph{tür}, o önermenin ispatı ise o türden bir \emph{program}
olarak görülebilir. ``$P$ ise $Q$'' önermesinin bir ispatı, $P$'nin
kanıtını alıp $Q$'nun kanıtını üreten bir fonksiyondur. Bu bakış açısı ilk
bölümlerde sezgisel olarak ortaya çıkar, Kısım~IV'te ise kesin bir matematiksel
teoreme dönüşür.

İkinci omurga, \textbf{yapıcı ve klasik mantık} arasındaki ayrımdır. \Lean{}
varsayılan olarak yapıcıdır: Üçüncü hâlin olanaksızlığını kullanmak
istediğinizde bunu açıkça söylemeniz gerekir. Bu durum, klasik ders
kitaplarında gözden kaçan bir soruyu okurun önüne koyar: \emph{Bir şeyin
yanlış olmadığını göstermek, doğru olduğunu göstermek midir?}

\section*{Eşlik eden Lean projesi}

Kitaptaki \Lean{} kodlarının tamamı, kitapla birlikte dağıtılan
\lc{ModernMantik} projesinden alınır. Kitabın \LaTeX{} kaynağı kodu bu
dosyalardan doğrudan okur; böylece kitapta basılı olan kod ile derlenen kod
her zaman aynıdır. Proje, otomatik bir sürekli tümleştirme (CI) iş akışıyla
her değişiklikte yeniden derlenir.

\vspace{1em}
\begin{flushright}
\yazar\\
\textit{Ekim 2026}
\end{flushright}
